\begin{description}
\item[(a)] From Figure 2, % sic: the period is read off the light curve, i.e. Figure 3 of the statement
  the period of the cepheid is $P \sim 11$ days and
  its average apparent magnitude is $\sim(14.8 + 14.1)/2$ mag, i.e.\
  $m = 14.45$ mag. \hfill(2 points)

  [A careful student will notice that the graph is not upside-down symmetrical
  so he/she choose some value closer to the bottom; that is $m = 14.5$ mag.
  \hfill(1 point)]

  From Figure 1, % sic: the period-luminosity relation is Figure 2 of the statement
  we derive that for a period of 11 days the expected absolute
  magnitude of the cepheid is $M \approx -4.2$. \hfill(1 point)

  [A careful student will notice that the graph is logarithmic so he/she
  choose a value closer to 4.3] \hfill(1 point)

  Using the formula $m - M = -5 + 5\log r$, where $r$ is the distance of the
  cepheid, we get
  \[ \log r = (14.45 + 4.2 + 5)/5 = 4.73, \]
  thence $r = 10^{4.73} \approx 57500$\,pc or $r = 57.5$\,kpc.
  \hfill(3 points)
\item[(b)] Assuming $A = 0.25$, then
  \[ \log r = (14.45 + 4.2 + 5 + 0.25)/5 = 4.78, \]
  % sic: the extinction term should be subtracted (m - M - A = -5 + 5 log r),
  % which gives log r = 4.68 and r ~ 48 kpc.
  thence $r = 10^{4.78} \sim 53000$\,pc or $r = 64.5$\,kpc.
  % sic: 10^{4.78} = 60300 pc = 60.3 kpc; neither printed value matches, and
  % 53000 pc is inconsistent with 64.5 kpc.
  \hfill(2 points)
\end{description}
